%==================================================================
% LAMPIRAN
%
% Setiap lampiran didaftarkan ke DAFTAR LAMPIRAN dengan perintah
% \daftarlampiran{...}
%
% Penomoran gambar dan tabel di dalam lampiran otomatis berubah
% menjadi A.1, A.2, B.1, dan seterusnya.
%==================================================================
\clearpage
\appendix
\renewcommand{\thefigure}{\Alph{chapter}.\arabic{figure}}
\renewcommand{\thetable}{\Alph{chapter}.\arabic{table}}

%------------------------------------------------------------------
\chapter{INSTRUMEN PENELITIAN}
\daftarlampiran{Lampiran A \quad Instrumen Penelitian}

Lampirkan instrumen lengkap yang digunakan, misalnya kuesioner, pedoman
wawancara, atau lembar validasi pakar.

\begin{table}[H]
    \centering
    \caption{Contoh butir instrumen}
    \begin{tabular}{|c|p{8cm}|c|}
        \hline
        \textbf{No} & \textbf{Pernyataan} & \textbf{Dimensi} \\ \hline
        1 & Saya merasa sistem mudah dipelajari.       & \textit{Usability} \\ \hline
        2 & Informasi yang disajikan sistem akurat.    & \textit{Information quality} \\ \hline
        3 & Sistem memberikan tanggapan yang memadai.  & \textit{Service interaction} \\ \hline
    \end{tabular}
\end{table}

%------------------------------------------------------------------
\chapter{DATA HASIL PENELITIAN}
\daftarlampiran{Lampiran B \quad Data Hasil Penelitian}

Lampirkan rekapitulasi data mentah, hasil perhitungan, atau keluaran perangkat
lunak analisis.

\begin{figure}[H]
    \centering
    \includegraphics[width=0.8\linewidth]{contoh-tangkapan-a}
    \caption{Contoh rekapitulasi data}
\end{figure}

%------------------------------------------------------------------
\chapter{DOKUMENTASI}
\daftarlampiran{Lampiran C \quad Dokumentasi}

Lampirkan dokumentasi kegiatan penelitian, surat izin, atau berita acara
pengujian.
